\section{Buyer-choice thresholds and comparison limits}
\label{app:buyer-thresholds}

This post-freeze appendix derives the frozen buyer-choice equations; it does not
run alternative worlds or revise their inputs. Its source is the registered
\texttt{economy.py} functions \texttt{utilities}, \texttt{opportunity}, and
\texttt{quoted\_precision}, with the frozen protocol. Exact source identities and
the longer derivation are retained in \texttt{BUYER\_THRESHOLDS.md}.

\textbf{Primary model limitation: asymmetric quality screening.}
The network must meet the buyer's explicit precision threshold, whereas the human
alternative is not subjected to that same hard gate. Human quality instead enters
utility as the fixed weight $q_H=0.985$. Thus a human offer may remain available
even when the buyer's threshold exceeds $0.985$. This is a stipulated comparison
asymmetry, not evidence that equivalent customers actually screen the two services
this way. Human delivery is also fixed at 20 model calendar days, a separate
assumption that can favor the network through discounting and deadline eligibility.

Let $V$ denote assumed project value, $P_A$ the network price, $P_H$ the hypothetical
human-service price, and $t_A,t_H$ their quoted delays in days. Let $q_A$ be the
network's public precision estimate conditional on technical acceptance, $d=0.85$
its fixed delivery weight, and $\tau$ the buyer's distinct trust discount. Customer
intake effort has dollar value $E$ on both sides; $S$ is the additional first-use
network switching burden. Write
$D_A=e^{-\delta t_A/7}$, $D_H=e^{-\delta t_H/7}$, and $w_A=\tau d q_A$, with weekly
decay $\delta=0.015$. Before eligibility restrictions, the exact utilities are
\begin{align}
 U_A &= w_A(VD_A-P_A)-E-S, \label{eq:buyer-network}\\
 U_H &= q_H(VD_H-P_H)-E,\\
 U_W &= \max\{0,\gamma q_H(Ve^{-\delta T_W}-P_H)-E-C_W\},
\end{align}
where postponement uses $\gamma=0.60$, $T_W=8$ weeks and $C_W=200$ USD; decline has
utility zero. These weights are assumptions, not measured model or professional
performance. Postponement is a reduced-form outside option, not a deferred-job
queue. Later opportunities are generated separately.

The network also requires $q_A\ge q_{\min}$, $P_A\le R$ and $t_A\le D$, where $R$
is the buyer's remaining annual budget and $D$ its deadline. The provider must be
open and able to fund and schedule the public maximum delivery reservation. The
human offer requires $P_H\le R$ and $t_H\le D$. Ineligible offers have utility
$-\infty$. Define $U_* = \max\{0,U_W,U_H\text{ if eligible}\}$. The network loses
ties to the earlier-listed outside options. Consequently, for $w_A>0$, its exact
one-opportunity adoption condition is
\begin{equation}
 \boxed{P_A < VD_A-\frac{E+S+U_*}{w_A}},
 \label{eq:buyer-complete-ceiling}
\end{equation}
in addition to the eligibility conditions. A nonpositive ceiling admits no
nonnegative winning price. If $w_A=0$, nonnegative effort and switching burdens
prevent the network from beating decline.

Against an eligible human offer alone, the ceiling simplifies to
\begin{equation}
 P_A < VD_A-\frac{q_H}{w_A}(VD_H-P_H)-\frac{S}{w_A}.
 \label{eq:buyer-human-ceiling}
\end{equation}
Common intake effort cancels in this pairwise comparison but still matters
against decline. Equivalently, when $VD_A>P_A$, required trust is
$\tau>(E+S+U_*)/[d q_A(VD_A-P_A)]$; a requirement above one is infeasible. These
analytical ceilings are not prices or private buyer information supplied to the
provider. Its adaptive controller learns from its own quote/conversion history.

\textbf{Human latency sensitivity.}
The fixed $t_H=20$ is neither observed professional turnaround nor inferred from
the 64/180/488-hour work breakdown; staffing, parallel work and expedited human
service are not modeled. Holding one buyer state fixed,
\begin{equation}
 \frac{\partial P_A^{\max,H}}{\partial t_H}
 =\frac{q_H}{w_A}\frac{\delta}{7}V e^{-\delta t_H/7}>0.
\end{equation}
Faster human service therefore lowers the human-relative network price ceiling.
It can also restore an otherwise excluded human offer: $D<20$ removes that offer
in the frozen model, affecting one fifth of the assumed uniform 10--60-day
deadline draws. This is not an empirical market share or an adoption estimate.

At a fixed state, applying the same hard precision gate to both services could
only remove human options; it cannot make the network less attractive. The frozen
AI-only gate therefore weakens, or leaves unchanged, the network's adoption case
relative to that symmetric counterfactual. This directional observation does not
establish an overall conservative model: the 20-day human-delay assumption and
other inputs can act in the opposite direction. Changes to choices would alter
later cash, queues, trust, learning and prices, so no dynamic-profit or investment
claim follows from these static inequalities. Profitable accepted delivery and
incremental retained Alpha remain separate tests.
